\section{有限閉被覆から極大な相異なる閉像へ}\label{sec:maximal-image-synthesis}
ここで全源の認識と成分の極大性の論理をまとめる。
有限閉像の一覧を先に答えと仮定せず、各全sourceのJacobi恒等式・逆写像・全補集合を
各入口で証明してから、この節を適用する。

\begin{lemma}[有限閉被覆の極大像]
\label{lem:finite-maximal-images}
有限型の被約空間$X$が空でない閉既約集合$C_1,\ldots,C_n$で覆われているとする。
全ての$i\ne j$について$C_i\not\subset C_j$なら、これらはちょうど$X$の全既約成分である。
また既に全空間成分と証明された$K_1,\ldots,K_m$と、次元$d$の既約閉像
$Y_1,\ldots,Y_r$で$X$を覆い、各$Y_a$が全$K_b$の外に出て、$Y_a$同士が相異なるなら、
各$Y_a$は全空間成分となる。
\end{lemma}
\begin{proof}
既約集合が有限閉和に入れば一つの項に含まれる。
従って各既約成分はある$C_i$に含まれ、逆の包含から一致する。
後半では$Y_a$を含む成分を選ぶ。それは既存$K_b$には含まれないので、ある$Y_c$に含まれる。
同次元の閉既約包含$Y_a\subset Y_c$は等号であり、相違から$a=c$となる。
\end{proof}

\subsection{本分類で用いる三つの成分性の証明}
第一は全周囲のJacobi微分による挟み撃ちである。
全150係数の核にはEuler鉛直24方向があり、非零coneの射影化でさらに一次元を除く。
実際のsource・群作用・尺度を含む微分の階数を$s$、全Jacobi微分の階数を$r$としたとき、
\[
 s-25\le\dim C_i\le\dim_p\mathcal P\le125-r.
\]
両端が一致すれば、実際の既約像は局所成分であり、その閉包は全空間成分である。
単に源の変数を数えた値や、補助条件を課したsliceの接空間を上界には使わない。

第二は補題\ref{lem:transverse-support}の真正形式横断切片である。
実sourceの方向を固定して線形可逆方向を解いた後の法線変数に、零次元の初期イデアルを与える。
線形解の任意性を戻した高次補正まで含めて初期式を証明するので、
有限jetが持ち上がるとの仮定も、schemeが一般に被約との仮定も要らない。
NEW25と初期の非被約成分、rank一のB1/B3/V1等はこの方法を用いる。
この場合の正確な入力・障害functional・行列・minorは対応命題の証明書を用いる。

第三は有限閉被覆による極大性である。
低modular旗の外での全像次元上界を証明してM1/M11を扱い、
最後の五つの23次元像には命題\ref{prop:five-last-full-components}と
補題\ref{lem:finite-maximal-images}を用いる。
そこでは全既存142成分との非包含を、全閉包の因子incidence、canonical cofactorの旗、
T36/T38の全proper flag、T37の有界次数葉曲面、実磁気商の不変量から証明した。
一点での接空間値や一般因子型が違うという事実だけではこの結論を出していない。
同じ23次元の最後の五像の相違も、本文の全磁気shearとmarked曲線の比較による。

\subsection{標識忘却・特殊底・階数二境界}
marked sourceを忘れる前に像の既約性を保ち、一般幾何学的fiberと有限な標識の置換を区別する。
DPMの二つの平方根、MPSの頂点置換、同じ自由平面曲線の異なるSaito frameから
新しい成分を加算しない。これらの実座標同値は各sourceの同定で処理した。
被約因子では定理\ref{thm:reduced-fibre-linear}の全fixed-D fiberが空でない線形開であるから、
系\ref{cor:reduced-divisor-component-images}により別成分の閉因子像は一致しない。
しかし因子像の包含だけから成分包含は結論せず、kernel rankが跳ぶ特殊底は各全支持証明へ戻す。
非被約因子にはこの線形化を流用せず、全Jacobiと高次補正を残す。

各A/B/C入口の全点被覆を実Poisson変形で先に閉じると、その有限閉和は
$\overline{\mathcal P^{(4)}}$も覆う。
その後で命題\ref{prop:rank2-support}の閉包項を代入し、六つの純rank二像を加える。
六像の全周囲成分性と相違は第\ref{sec:rank-two-reconstruction}節による。
rank四成分とrank二成分は一般rankが異なり、一致しない。
以上の順序はrank二の支持定理にrank四分類を循環的に仮定するものではない。

このように得た相異なる極大閉像を第\ref{sec:component-catalogue}節で数えると、
rank四では$40+11+1+89=141$、rank二では六となる。
次元別個数は同節の各既約源の証明済み像次元を集計したものであり、
場合分けを打ち切る公理として用いた数ではない。
